\chapter{Reference Implementations}
\label{app:implementations}

This appendix provides self-contained Python examples that illustrate the
algorithms discussed throughout the book. The educational LWE example uses
only standard libraries; all remaining examples require \texttt{liboqs-python}
(\url{https://github.com/open-quantum-safe/liboqs-python}).
Install it with \texttt{pip install liboqs-python}.

\medskip
\noindent\textbf{Warning.} None of these scripts are suitable for production
use. They are intended solely as pedagogical companions to the main text.

%% ==============================================================
\section{LWE Encryption (Educational)}
\label{app:lwe-impl}

The following pure-Python script implements a toy LWE encryption scheme with
parameters $n=16$, $q=97$. These parameters provide \emph{no} security;
they are chosen so that every intermediate value fits comfortably on screen.
The code demonstrates the three core operations---key generation, single-bit
encryption, and decryption---and highlights how noise accumulation can cause
decryption failures when parameters are too aggressive.

\begin{lstlisting}[language=Python,caption={Toy LWE encryption (NOT secure)}]
import random

# --- Parameters (educational only) ---
n = 16    # lattice dimension
q = 97    # modulus (prime)
m = 32    # number of public-key samples
sigma = 2 # noise bound (uniform in [-sigma, sigma])

def sample_error():
    """Sample small noise uniformly from {-sigma, ..., sigma}."""
    return random.randint(-sigma, sigma) % q

def keygen():
    """Return (public_key, secret_key)."""
    s = [random.randrange(q) for _ in range(n)]        # secret vector
    A = [[random.randrange(q) for _ in range(n)]
         for _ in range(m)]                             # random matrix
    e = [sample_error() for _ in range(m)]              # noise vector
    b = [(sum(A[i][j]*s[j] for j in range(n)) + e[i]) % q
         for i in range(m)]                             # b = As + e
    return (A, b), s

def encrypt(pk, bit):
    """Encrypt a single bit in {0, 1}."""
    A, b = pk
    # Choose a random binary selector vector r
    r = [random.randint(0, 1) for _ in range(m)]
    # Sum selected rows of A and corresponding entries of b
    u = [sum(r[i]*A[i][j] for i in range(m)) % q
         for j in range(n)]
    v = sum(r[i]*b[i] for i in range(m)) % q
    # Encode the bit: add floor(q/2) if bit == 1
    v = (v + bit * (q // 2)) % q
    return (u, v)

def decrypt(sk, ct):
    """Decrypt ciphertext to recover the bit."""
    u, v = ct
    s = sk
    # Compute v - <u, s> mod q
    d = (v - sum(u[j]*s[j] for j in range(n))) % q
    # Map to the nearest bit: 0 if close to 0, 1 if close to q/2
    if d > q // 2:
        d -= q  # centered representative
    return 0 if abs(d) < q // 4 else 1

# --- Quick self-test ---
pk, sk = keygen()
for bit in (0, 1):
    ct = encrypt(pk, bit)
    assert decrypt(sk, ct) == bit, "Decryption failure!"
print("LWE encrypt/decrypt OK for bits 0 and 1")
\end{lstlisting}

%% ==============================================================
\section{ML-KEM with liboqs}
\label{app:mlkem-impl}

ML-KEM (FIPS~203) is the NIST standard key-encapsulation mechanism derived
from CRYSTALS-Kyber. The script below exercises all three parameter sets
through the \texttt{liboqs} Python bindings, prints object sizes, and
measures wall-clock time for each operation.

\begin{lstlisting}[language=Python,caption={ML-KEM key encapsulation with liboqs}]
import oqs
import time

PARAM_SETS = ["ML-KEM-512", "ML-KEM-768", "ML-KEM-1024"]
ITERS = 100

for name in PARAM_SETS:
    kem = oqs.KeyEncapsulation(name)

    # --- Key generation ---
    t0 = time.perf_counter()
    for _ in range(ITERS):
        pk = kem.generate_keypair()
    kg_ms = (time.perf_counter() - t0) / ITERS * 1000

    # --- Encapsulation ---
    t0 = time.perf_counter()
    for _ in range(ITERS):
        ct, ss_enc = kem.encap_secret(pk)
    enc_ms = (time.perf_counter() - t0) / ITERS * 1000

    # --- Decapsulation ---
    t0 = time.perf_counter()
    for _ in range(ITERS):
        ss_dec = kem.decap_secret(ct)
    dec_ms = (time.perf_counter() - t0) / ITERS * 1000

    assert ss_enc == ss_dec, "Shared secrets do not match!"

    print(f"\n=== {name} ===")
    print(f"  Public key : {len(pk):>5} bytes")
    print(f"  Ciphertext : {len(ct):>5} bytes")
    print(f"  Shared sec.: {len(ss_enc):>5} bytes")
    print(f"  KeyGen     : {kg_ms:>7.3f} ms")
    print(f"  Encaps     : {enc_ms:>7.3f} ms")
    print(f"  Decaps     : {dec_ms:>7.3f} ms")
\end{lstlisting}

%% ==============================================================
\section{ML-DSA with liboqs}
\label{app:mldsa-impl}

ML-DSA (FIPS~204), derived from CRYSTALS-Dilithium, is the primary NIST
post-quantum signature standard. The code below demonstrates key generation,
signing, and verification for ML-DSA-65 (the recommended security level)
and reports sizes and timings.

\begin{lstlisting}[language=Python,caption={ML-DSA digital signatures with liboqs}]
import oqs
import time

sig = oqs.Signature("ML-DSA-65")
message = b"Post-quantum cryptography is here."
ITERS = 100

# --- Key generation ---
t0 = time.perf_counter()
for _ in range(ITERS):
    pk = sig.generate_keypair()
kg_ms = (time.perf_counter() - t0) / ITERS * 1000

# --- Signing ---
t0 = time.perf_counter()
for _ in range(ITERS):
    signature = sig.sign(message)
sign_ms = (time.perf_counter() - t0) / ITERS * 1000

# --- Verification ---
verifier = oqs.Signature("ML-DSA-65")
t0 = time.perf_counter()
for _ in range(ITERS):
    valid = verifier.verify(message, signature, pk)
ver_ms = (time.perf_counter() - t0) / ITERS * 1000

assert valid, "Signature verification failed!"

print("=== ML-DSA-65 ===")
print(f"  Public key : {len(pk):>5} bytes")
print(f"  Signature  : {len(signature):>5} bytes")
print(f"  KeyGen     : {kg_ms:>7.3f} ms")
print(f"  Sign       : {sign_ms:>7.3f} ms")
print(f"  Verify     : {ver_ms:>7.3f} ms")
\end{lstlisting}

%% ==============================================================
\section{SLH-DSA with liboqs}
\label{app:slhdsa-impl}

SLH-DSA (FIPS~205), based on SPHINCS$^+$, is a stateless hash-based
signature scheme. Its main trade-off is between the ``small'' variants (shorter
signatures, slower signing) and the ``fast'' variants (larger signatures,
faster signing). The script below makes this trade-off concrete.

\begin{lstlisting}[language=Python,caption={SLH-DSA: small vs.\ fast parameter comparison}]
import oqs
import time

VARIANTS = [
    ("SLH-DSA-SHA2-128s", "128-bit, small"),
    ("SLH-DSA-SHA2-128f", "128-bit, fast"),
]
message = b"Hash-based signatures need no lattice assumptions."
ITERS = 10  # SLH-DSA signing is slow; fewer iterations

for name, label in VARIANTS:
    sig = oqs.Signature(name)
    pk = sig.generate_keypair()

    # --- Signing ---
    t0 = time.perf_counter()
    for _ in range(ITERS):
        signature = sig.sign(message)
    sign_ms = (time.perf_counter() - t0) / ITERS * 1000

    # --- Verification ---
    verifier = oqs.Signature(name)
    t0 = time.perf_counter()
    for _ in range(ITERS):
        valid = verifier.verify(message, signature, pk)
    ver_ms = (time.perf_counter() - t0) / ITERS * 1000

    assert valid
    print(f"\n=== {name} ({label}) ===")
    print(f"  Public key : {len(pk):>6} bytes")
    print(f"  Signature  : {len(signature):>6} bytes")
    print(f"  Sign       : {sign_ms:>9.1f} ms")
    print(f"  Verify     : {ver_ms:>9.1f} ms")
\end{lstlisting}

%% ==============================================================
\section{Hybrid TLS with Post-Quantum KEM}
\label{app:hybrid-tls}

Deploying post-quantum algorithms in TLS\,1.3 requires an OpenSSL build
that supports the OQS provider. The following commands assume OpenSSL~3.x
compiled with the \texttt{oqs-provider} plugin (see
\url{https://github.com/open-quantum-safe/oqs-provider} for build
instructions).

\medskip
\noindent\textbf{Step 1. Generate a hybrid certificate.}
The example below creates an ML-DSA-65 CA and issues a server certificate
that uses a hybrid key exchange during the TLS handshake.

\begin{lstlisting}[language=bash,caption={Generate an ML-DSA-65 CA and server certificate}]
# Generate CA key and self-signed certificate
openssl req -x509 -new -newkey mldsa65 -keyout ca.key \
  -out ca.crt -nodes -subj "/CN=PQC Test CA" -days 365

# Generate server key and CSR
openssl req -new -newkey mldsa65 -keyout server.key \
  -out server.csr -nodes -subj "/CN=localhost"

# Sign the server certificate with the CA
openssl x509 -req -in server.csr -CA ca.crt -CAkey ca.key \
  -CAcreateserial -out server.crt -days 365
\end{lstlisting}

\noindent\textbf{Step 2. Test a hybrid TLS connection.}
Use \texttt{s\_server} and \texttt{s\_client} to verify the handshake
with a hybrid key exchange group (X25519 + ML-KEM-768).

\begin{lstlisting}[language=bash,caption={Hybrid TLS 1.3 handshake test}]
# Terminal 1 -- start the server
openssl s_server -cert server.crt -key server.key \
  -groups x25519_mlkem768 -www -accept 4433

# Terminal 2 -- connect as a client
openssl s_client -connect localhost:4433 \
  -groups x25519_mlkem768 -CAfile ca.crt
\end{lstlisting}

\noindent A successful connection will display the negotiated group
(\texttt{x25519\_mlkem768}) in the handshake output, confirming that
both classical and post-quantum key exchange took place.

%% ==============================================================
\section{Performance Benchmarks}
\label{app:benchmarks}

The script below benchmarks every NIST-standardised PQC algorithm alongside
RSA-2048 and ECDSA-P256, printing the results as a formatted table. It
requires \texttt{liboqs-python} and the \texttt{cryptography} package.

\begin{lstlisting}[language=Python,caption={Comprehensive PQC benchmark script}]
import oqs, time
from cryptography.hazmat.primitives.asymmetric import rsa, ec, padding, utils
from cryptography.hazmat.primitives import hashes, serialization

ITERS = 100
MSG = b"Benchmark payload for signature schemes."

def bench(fn, n=ITERS):
    t0 = time.perf_counter()
    for _ in range(n):
        result = fn()
    return (time.perf_counter() - t0) / n * 1000, result

rows = []

# --- ML-KEM ---
for name in ("ML-KEM-512", "ML-KEM-768", "ML-KEM-1024"):
    k = oqs.KeyEncapsulation(name)
    kg, pk = bench(k.generate_keypair)
    en, (ct, ss) = bench(lambda: k.encap_secret(pk))
    de, _  = bench(lambda: k.decap_secret(ct))
    rows.append((name, f"{len(pk)}", f"{len(ct)}",
                 f"{kg:.3f}", f"{en:.3f}", f"{de:.3f}"))

# --- ML-DSA ---
for name in ("ML-DSA-44", "ML-DSA-65", "ML-DSA-87"):
    s = oqs.Signature(name)
    kg, pk = bench(s.generate_keypair)
    si, sig = bench(lambda: s.sign(MSG))
    v = oqs.Signature(name)
    ve, _ = bench(lambda: v.verify(MSG, sig, pk))
    rows.append((name, f"{len(pk)}", f"{len(sig)}",
                 f"{kg:.3f}", f"{si:.3f}", f"{ve:.3f}"))

# --- SLH-DSA (fewer iterations -- signing is slow) ---
for name in ("SLH-DSA-SHA2-128s", "SLH-DSA-SHA2-128f"):
    s = oqs.Signature(name)
    kg, pk = bench(s.generate_keypair, 10)
    si, sig = bench(lambda: s.sign(MSG), 10)
    v = oqs.Signature(name)
    ve, _ = bench(lambda: v.verify(MSG, sig, pk), 10)
    rows.append((name, f"{len(pk)}", f"{len(sig)}",
                 f"{kg:.3f}", f"{si:.3f}", f"{ve:.3f}"))

# --- RSA-2048 (sign/verify only) ---
sk = rsa.generate_private_key(65537, 2048)
pk_rsa = sk.public_key()
sig_rsa = sk.sign(MSG, padding.PKCS1v15(), hashes.SHA256())
kg_r, _ = bench(lambda: rsa.generate_private_key(65537, 2048), 10)
si_r, _ = bench(lambda: sk.sign(MSG, padding.PKCS1v15(),
                                hashes.SHA256()))
ve_r, _ = bench(lambda: pk_rsa.verify(sig_rsa, MSG,
                                      padding.PKCS1v15(),
                                      hashes.SHA256()))
pk_bytes = pk_rsa.public_bytes(serialization.Encoding.DER,
             serialization.PublicFormat.SubjectPublicKeyInfo)
rows.append(("RSA-2048", f"{len(pk_bytes)}", f"{len(sig_rsa)}",
             f"{kg_r:.3f}", f"{si_r:.3f}", f"{ve_r:.3f}"))

# --- ECDSA P-256 ---
sk_ec = ec.generate_private_key(ec.SECP256R1())
pk_ec = sk_ec.public_key()
sig_ec = sk_ec.sign(MSG, ec.ECDSA(hashes.SHA256()))
kg_e, _ = bench(lambda: ec.generate_private_key(ec.SECP256R1()))
si_e, _ = bench(lambda: sk_ec.sign(MSG, ec.ECDSA(hashes.SHA256())))
ve_e, _ = bench(lambda: pk_ec.verify(sig_ec, MSG,
                                     ec.ECDSA(hashes.SHA256())))
pk_ec_b = pk_ec.public_bytes(serialization.Encoding.DER,
            serialization.PublicFormat.SubjectPublicKeyInfo)
rows.append(("ECDSA-P256", f"{len(pk_ec_b)}", f"{len(sig_ec)}",
             f"{kg_e:.3f}", f"{si_e:.3f}", f"{ve_e:.3f}"))

# --- Print table ---
hdr = f"{'Algorithm':<20} {'PK(B)':>7} {'CT/Sig(B)':>9} " \
      f"{'KG(ms)':>8} {'Op1(ms)':>8} {'Op2(ms)':>8}"
print(hdr)
print("-" * len(hdr))
for r in rows:
    print(f"{r[0]:<20} {r[1]:>7} {r[2]:>9} "
          f"{r[3]:>8} {r[4]:>8} {r[5]:>8}")
\end{lstlisting}

\noindent In the table above, \textbf{Op1} denotes encapsulation (for KEMs)
or signing (for signatures), and \textbf{Op2} denotes decapsulation or
verification, respectively. Readers are encouraged to run the script on
their own hardware and compare results with the figures in the main text.
